\documentclass[aspectratio=169]{beamer}

% Compile with pdfLaTeX.
\usepackage[T2A]{fontenc}
\usepackage[utf8]{inputenc}
\usepackage[russian,english]{babel}
\usepackage{graphicx}
\usepackage[defaultsans]{opensans}
\usepackage{tabularx}
\usefonttheme{professionalfonts}
\renewcommand{\familydefault}{\sfdefault}

\setbeamercolor{normal text}{fg=black}
\setbeamercolor{structure}{fg=black}
\setbeamercolor{frametitle}{fg=black}
\setbeamercolor{footline}{fg=black}
\setbeamercolor{item}{fg=black}
\setbeamercolor{subitem}{fg=black}
\setbeamercolor{subsubitem}{fg=black}
\setbeamercolor{title}{fg=white}

\setbeamerfont{normal text}{family*=\opensansfamily}
\setbeamerfont{title}{family*=\opensansfamily,size=\Huge,series=\bfseries}
\setbeamerfont{frametitle}{family*=\opensansfamily,series=\bfseries}
\setbeamerfont{footline}{family*=\opensansfamily}
\setbeamerfont{author}{family*=\opensansfamily}
\setbeamerfont{date}{family*=\opensansfamily}
\setbeamerfont{institute}{family*=\opensansfamily}

\setbeamertemplate{navigation symbols}{}
\setbeamertemplate{footline}[frame number]
\setbeamertemplate{frametitle}{%
  \vspace*{0.5cm}% Move titles slightly down
  \nointerlineskip
  \begin{beamercolorbox}[wd=\paperwidth,ht=2.8ex,dp=1.1ex,leftskip=5.7cm]{frametitle}
    \usebeamerfont{frametitle}\insertframetitle
  \end{beamercolorbox}%
}

\vspace*{1cm}
\title{Распознавание рукописного текста с анализом эмоциональной окраски}
\date{}

\newcommand{\slidebg}[1]{%
  \setbeamertemplate{background}{%
    \parbox[c][\paperheight][c]{\paperwidth}{%
      \includegraphics[width=\paperwidth,height=\paperheight]{#1}%
    }%
  }%
}

\begin{document}


\slidebg{assets/1.png}
\begin{frame}[plain,noframenumbering]
  \titlepage
\end{frame}


\slidebg{assets/2.png}
\begin{frame}
  \frametitle{Актуальность}
  \begin{itemize}
    \item Цифровизация рукописных архивов и исторических документов.
    \item Снижение ручного труда при вводе текстов из бумажных источников.
    \item Сокращение числа ошибок и времени обработки данных.
    \item Применение в образовании, медицине и документообороте.
    \item Основа для последующего анализа содержания и тональности текста.
  \end{itemize}
\end{frame}

\slidebg{assets/2.png}
\begin{frame}
  \frametitle{Цель и задачи работы}
  \scriptsize
  \textbf{Цель работы:} разработать программный инструмент для автоматического распознавания рукописного текста и последующего анализа тональности распознанного содержания.

  \vspace{0.3em}
  \textbf{Задачи работы:}
  \begin{columns}[T,onlytextwidth]
    \begin{column}{0.49\textwidth}
      \begin{itemize}
        \setlength{\itemsep}{0.2em}
        \item Проанализировать существующие подходы к распознаванию рукописного текста и анализу тональности.
        \item Подобрать и подготовить датасет для обучения и тестирования модели.
        \item Определить метрики качества и критерии успешности решения.
        \item Разработать и реализовать модель распознавания рукописного текста.
      \end{itemize}
    \end{column}
    \begin{column}{0.49\textwidth}
      \begin{itemize}
        \setlength{\itemsep}{0.2em}
        \item Обучить модель и оценить качество её работы на выбранном датасете.
        \item Реализовать модуль анализа тональности распознанного текста.
        \item Объединить распознавание и анализ тональности в единый программный инструмент.
      \end{itemize}
    \end{column}
  \end{columns}
\end{frame}

\slidebg{assets/2.png}
\begin{frame}
  \frametitle{Подходы к распознаванию: плюсы и минусы}
  \scriptsize
  \setlength{\tabcolsep}{4pt}
  \renewcommand{\arraystretch}{1.15}
  \begin{tabularx}{\textwidth}{|X|X|X|X|}
    \hline
    \textbf{Подход} & \textbf{Плюсы} & \textbf{Минусы} & \textbf{Где лучше работает} \\
    \hline
    Классический OCR (правила, сегментация) & Простота, низкие требования к ресурсам & Плохо переносит сложный почерк и шум & Печатный или аккуратный текст \\
    \hline
    CRNN + CTC & Не нужна посимвольная разметка, устойчив к разной длине строк, сильный базовый результат & Чувствителен к качеству предобработки, требует обучения & Рукописные строки средней сложности \\
    \hline
    Transformer-подходы (TrOCR и аналоги) & Высокий потенциал качества, гибкость архитектуры & Высокая вычислительная стоимость, сложнее обучение & Большие датасеты и мощные вычисления \\
    \hline
  \end{tabularx}

  \vfill
  \textbf{Анализ показал, что наиболее результативным является метод CRNN + CTC.}
\end{frame}

% \slidebg{assets/2.png}
% \begin{frame}
%   \frametitle{План презентации}
%   \begin{enumerate}
%     \item Постановка задачи и актуальность.
%     \item Обзор существующих подходов и их сравнение.
%     \item Обоснование выбора метода CRNN + CTC.
%     \item Описание нашего решения и пайплайна обработки.
%     \item Датасет IAM, метрики CER/WER и протокол оценки.
%     \item Текущий статус и дальнейшие шаги проекта.
%   \end{enumerate}
% \end{frame}

% \slidebg{assets/2.png}
% \begin{frame}
%   \frametitle{Наш метод: общая схема}
%   \begin{itemize}
%     \item \textbf{Вход:} изображение рукописного текста.
%     \item \textbf{Предобработка:} бинаризация, выравнивание строк, ресайз, нормализация.
%     \item \textbf{Модель:} CRNN извлекает признаки и моделирует последовательность.
%     \item \textbf{Декодирование:} CTC-декодер преобразует вероятности в текст.
%     \item \textbf{Дополнительно:} лексиконный анализ эмоциональной окраски распознанного текста.
%   \end{itemize}
% \end{frame}

% \slidebg{assets/2.png}
% \begin{frame}
%   \frametitle{Наш метод: предобработка данных}
%   \begin{itemize}
%     \item Удаление шумов и повышение контрастности.
%     \item Бинаризация изображения для выделения штрихов.
%     \item Геометрическое выравнивание строк текста.
%     \item Приведение к фиксированному размеру и нормализация.
%     \item Цель: стабилизировать вход для обучения и инференса CRNN.
%   \end{itemize}
% \end{frame}

% \slidebg{assets/2.png}
% \begin{frame}
%   \frametitle{Наш метод: модель CRNN + CTC}
%   \begin{itemize}
%     \item CNN-блок извлекает локальные визуальные признаки символов.
%     \item RNN-блок учитывает контекст по всей строке.
%     \item CTC-loss обучает модель без точного посимвольного выравнивания.
%     \item CER и WER используются как ключевые метрики качества.
%     \item Датасет: IAM Handwriting Database (строчная разметка и транскрипции).
%   \end{itemize}
% \end{frame}

% \slidebg{assets/3.png}
% \begin{frame}
%   \frametitle{Текущий статус и дальнейшие шаги}
%   \begin{itemize}
%     \item Реализуется базовый пайплайн распознавания рукописного текста.
%     \item Следующий этап: экспериментальная оценка на IAM (CER/WER).
%     \item Планируется сравнение с альтернативными архитектурами.
%     \item Побочная задача: интеграция анализа тональности распознанного текста.
%   \end{itemize}
% \end{frame}

% \slidebg{assets/3.png}
% \begin{frame}
%   \centering
%   \Huge Спасибо за внимание!
% \end{frame}

\end{document}
